% スイング・ツイスト分解
\begin{tikzpicture}[>=Stealth,font=\small]
\begin{scope}
  \fill (0,0) circle (2pt) node[left]{肩};
  \draw[line width=4pt,gray!50] (0,0)--(0,-2);
  \draw[line width=4pt,blue!60] (0,0)--(1.4,-1.4);
  \draw[->,thick] (0,-1.2) arc (-90:-45:1.2);
  \node at (0.75,-1.35) {スイング};
  \node[align=center,font=\footnotesize] at (0.6,-2.6) {腕の向きを変える\\（2 自由度）};
\end{scope}
\begin{scope}[shift={(4.5,0)}]
  \fill (0,0) circle (2pt);
  \draw[line width=4pt,blue!60] (0,0)--(1.4,-1.4);
  \draw[thick] (0.9,-0.9) ellipse (0.35 and 0.15);
  \draw[->,thick] (1.25,-0.9) arc (0:-200:0.35 and 0.15);
  \node at (1.8,-0.6) {ツイスト};
  \node[align=center,font=\footnotesize] at (0.6,-2.6) {腕の長軸まわりにひねる\\（1 自由度）};
\end{scope}
\node[font=\footnotesize,align=left] at (9.2,-1.0) {上腕の姿勢 $R = R_{\text{swing}}R_{\text{twist}}$\\ひねりの値は分解の仕方\\（基準姿勢・順序）に依存する};
\end{tikzpicture}
